% =============================================================================
% Data Science Theory: An Intuitive Guide
% Author: Andres Gonzalez Ortega
% =============================================================================
\documentclass[12pt,a4paper]{article}
\usepackage[utf8]{inputenc}
\usepackage[T1]{fontenc}
\usepackage{lmodern}
\usepackage{microtype}
\usepackage[margin=2.8cm, top=3cm, bottom=3cm]{geometry}
\usepackage{fancyhdr}
\usepackage{titlesec}
\usepackage{amsmath, amssymb}
\usepackage{mathtools}
\usepackage{bm}
\usepackage{booktabs}
\usepackage{enumitem}
\usepackage[most]{tcolorbox}
\usepackage{xcolor}
\usepackage{hyperref}
\usepackage{parskip}
\setlength{\headheight}{14pt}

% ---------------------------------------------------------------------------
% Colors
% ---------------------------------------------------------------------------
\definecolor{conceptpurple}{RGB}{110, 64, 170}
\definecolor{conceptpurplebg}{RGB}{243, 238, 252}
\definecolor{analogyamber}{RGB}{200, 130, 20}
\definecolor{analogyamberbg}{RGB}{254, 245, 224}
\definecolor{mentalgreen}{RGB}{34, 120, 92}
\definecolor{mentalgreenbg}{RGB}{232, 246, 240}
\definecolor{accentblue}{RGB}{28, 96, 185}
\definecolor{inkdark}{RGB}{32, 32, 40}
\definecolor{inkmid}{RGB}{96, 100, 112}

% ---------------------------------------------------------------------------
% Box environments
% ---------------------------------------------------------------------------
\newtcolorbox{concept}[1][Concept]{enhanced, breakable,
  colback=conceptpurplebg, colframe=conceptpurple,
  fonttitle=\bfseries\color{conceptpurple}, title={#1},
  titlerule=0pt, sharp corners=southwest, rounded corners=northeast,
  boxrule=0pt, leftrule=3pt, top=8pt, bottom=8pt, left=10pt, right=10pt}

\newtcolorbox{analogy}[1][Analogy]{enhanced, breakable,
  colback=analogyamberbg, colframe=analogyamber,
  fonttitle=\bfseries\color{analogyamber}, title={#1},
  titlerule=0pt, sharp corners=southwest, rounded corners=northeast,
  boxrule=0pt, leftrule=3pt, top=8pt, bottom=8pt, left=10pt, right=10pt}

\newtcolorbox{mentalmodel}[1][Mental Model]{enhanced, breakable,
  colback=mentalgreenbg, colframe=mentalgreen,
  fonttitle=\bfseries\color{mentalgreen}, title={#1},
  titlerule=0pt, sharp corners=southwest, rounded corners=northeast,
  boxrule=0pt, leftrule=3pt, top=8pt, bottom=8pt, left=10pt, right=10pt}

% ---------------------------------------------------------------------------
% Section title format
% ---------------------------------------------------------------------------
\titleformat{\section}{\Large\bfseries\color{inkdark}}{}{0pt}{}
\titleformat{\subsection}{\large\bfseries\color{inkmid}}{}{0pt}{}

% ---------------------------------------------------------------------------
% Header / footer
% ---------------------------------------------------------------------------
\pagestyle{fancy}
\fancyhf{}
\fancyhead[R]{\small\color{inkmid}\thepage}
\renewcommand{\headrulewidth}{0pt}
\renewcommand{\footrulewidth}{0pt}

% ---------------------------------------------------------------------------
% Hyperref
% ---------------------------------------------------------------------------
\hypersetup{colorlinks=true, linkcolor=accentblue, urlcolor=accentblue,
  pdfauthor={Andres Gonzalez Ortega},
  pdftitle={Data Science Theory: An Intuitive Guide}}

% =============================================================================
\begin{document}
% =============================================================================

% ---------------------------------------------------------------------------
% Title block
% ---------------------------------------------------------------------------
\thispagestyle{empty}
\vspace*{2cm}
\begin{flushleft}
  {\Huge\bfseries\color{inkdark} Data Science Theory}\\[10pt]
  {\large\color{inkmid} An intuitive guide from probability to production}\\[20pt]
  {\small\color{inkmid} Andres Gonzalez Ortega
    \quad\textbullet\quad \today}
\end{flushleft}
\vspace{1.5cm}

% ---------------------------------------------------------------------------
% Table of contents
% ---------------------------------------------------------------------------
\tableofcontents
\newpage

% ---------------------------------------------------------------------------
% Reading Guide
% ---------------------------------------------------------------------------
\section*{How to Read This Document}
\addcontentsline{toc}{section}{How to Read This Document}

This document uses three colored callout boxes. Each color carries a fixed
meaning, so the rest of the document leaves the boxes untitled --- the color
alone tells you which kind of content you are reading.

\begin{concept}[Concept --- purple]
The crisp definition of the idea. Answers \emph{what} the concept is.
\end{concept}

\begin{analogy}[Analogy --- amber]
A familiar situation that maps onto the concept. Answers what the concept
\emph{feels like}.
\end{analogy}

\begin{mentalmodel}[Mental Model --- green]
The picture to keep in mind when you recall the idea later. Answers how to
\emph{remember} it.
\end{mentalmodel}

\vspace{10pt}
\hrule height 0.4pt
\vspace{14pt}

% =============================================================================
\part*{Part I: Mathematical Foundations}
\addcontentsline{toc}{part}{Part I: Mathematical Foundations}
% =============================================================================

% ---------------------------------------------------------------------------
\section*{Probability and Statistical Estimation}
\addcontentsline{toc}{section}{Probability and Statistical Estimation}
% ---------------------------------------------------------------------------

Before fitting any model, it is worth asking: what does it mean for an
estimate to be good? Maximum likelihood estimation gives one precise answer,
and the answer has a beautiful geometric interpretation. How confident should
we be in any particular estimate, and where does that confidence actually
come from?

\begin{concept}
Maximum likelihood estimation chooses the parameter value that makes the
observed data most probable. Its key theoretical payoff is asymptotic
normality: as the sample grows, the MLE concentrates around the true value
at a rate governed by the Fisher information, and that concentration
justifies standard errors and Wald confidence intervals across virtually
all of applied statistics.
\end{concept}

\begin{analogy}
Think of the likelihood function as a landscape, and the MLE as the summit.
Fisher information measures how sharply peaked that summit is. A sharp peak
means the data pins down the parameter tightly; a broad, flat hill means
many parameter values fit almost equally well, so any single estimate is
uncertain. The variance of the MLE is literally the reciprocal of that
curvature.
\end{analogy}

\begin{equation*}
  \sqrt{n}\,\bigl(\hat\theta_{\mathrm{MLE}} - \theta_0\bigr)
  \xrightarrow{d}
  \mathcal{N}\!\left(0,\; I(\theta_0)^{-1}\right)
\end{equation*}

Here $I(\theta_0)$ is the Fisher information matrix evaluated at the true
parameter and $n$ is the sample size.

\begin{mentalmodel}
The bridge from MLE to actuarial credibility runs through the normal-normal
conjugate model. The Buhlmann credibility formula
$Z\bar{X}_i + (1-Z)\mu$ is precisely the Bayesian posterior mean; MLE
estimates the hyperparameters $(\mu, k)$ from the portfolio data.
Asymptotic normality is the frequentist road to the same destination.
\end{mentalmodel}

\vspace{10pt}\hrule height 0.4pt\vspace{14pt}

% ---------------------------------------------------------------------------
\section*{Optimization and Loss Functions}
\addcontentsline{toc}{section}{Optimization and Loss Functions}
% ---------------------------------------------------------------------------

Every machine learning model is, at its core, a minimization problem.
The choice of loss function decides what the model is actually optimized
to do; the choice of optimizer decides how efficiently it gets there.
Have you ever wondered why two models with identical architectures can
behave so differently in practice?

\begin{concept}
Every ML model is defined by
$\hat\theta = \arg\min_\theta\, L(\theta;\,\text{data})$.
The loss function and optimizer are not implementation details; they
determine what the model is solving and whether it will find a good
solution. Changing the loss can be more impactful than changing the
architecture.
\end{concept}

\begin{analogy}
Picture SGD momentum descending a narrow, steep-sided ravine. Without
momentum, the gradient oscillates back and forth across the high-curvature
walls and makes slow progress along the valley floor. With momentum, the
running average of gradients damps those side-to-side oscillations and
amplifies the consistent downhill direction, reaching the bottom far more
quickly.
\end{analogy}

The AdamW update combines adaptive learning rates with decoupled weight
decay:

\begin{equation*}
  \theta_{t+1} = \theta_t
  - \eta\!\left(\frac{\hat{m}_t}{\sqrt{\hat{v}_t} + \epsilon}
    + \lambda\,\theta_t\right)
\end{equation*}

where $\hat{m}_t$ and $\hat{v}_t$ are bias-corrected first and second
gradient moments and $\lambda\,\theta_t$ is the decoupled weight-decay
term that penalizes large weights independently of the gradient magnitude.

\begin{mentalmodel}
The condition number $\kappa$ of the loss surface is its aspect ratio:
a wide, shallow bowl in one direction and a steep wall in another.
Adaptive optimizers implicitly rescale each parameter's learning rate by
the inverse of its recent gradient magnitude, transforming that
ill-conditioned bowl into something closer to a sphere, where every
direction is equally easy to traverse.
\end{mentalmodel}

\vspace{10pt}\hrule height 0.4pt\vspace{14pt}

% ---------------------------------------------------------------------------
\section*{A Note on the Glossary}
\addcontentsline{toc}{section}{A Note on the Glossary}
% ---------------------------------------------------------------------------

Much confusion in data science comes from the same word meaning different
things in different communities. A ``model'' in statistics is a probability
distribution; in machine learning it is a function from inputs to
predictions; in business it is any analytical tool used to make a decision.
Which definition is in play matters more than we usually admit.

\begin{concept}
Precision in language is precision in thinking. Before presenting results
to a new audience, anchor every key term: what is the outcome variable,
what counts as a prediction, and what criterion defines success. A number
without a unit is meaningless; a metric without a definition is worse,
because it creates false confidence.
\end{concept}

\begin{analogy}
Consider the word ``significant.'' In statistics it means the p-value fell
below a pre-specified threshold; in business it means the effect is large
enough to act on; in casual speech it means noteworthy. The same word in
a board presentation and in a methods section is doing completely different
work.
\end{analogy}

\begin{mentalmodel}
Keep a one-page glossary for every project and share it at the first
stakeholder meeting. The questions it provokes are the most valuable
part of the exercise: they surface implicit disagreements about what is
actually being measured before anyone has spent a week modeling the wrong
thing.
\end{mentalmodel}

\vspace{10pt}\hrule height 0.4pt\vspace{14pt}

% =============================================================================
\part*{Part II: Generalization, Validation, and Feature Engineering}
\addcontentsline{toc}{part}{Part II: Generalization, Validation, and Feature Engineering}
% =============================================================================

% ---------------------------------------------------------------------------
\section*{Bias, Variance, and the Generalization Gap}
\addcontentsline{toc}{section}{Bias, Variance, and the Generalization Gap}
% ---------------------------------------------------------------------------

A model that memorizes the training data perfectly is not necessarily a
good model. The central question is always: how does it perform on data
it has never seen? The bias-variance decomposition makes that question
mathematically precise.

\begin{concept}
For squared-error loss, expected prediction error decomposes into three
parts: bias squared (systematic error that persists across different
training sets), variance (how much the model fluctuates as training data
changes), and irreducible noise from the world itself. Reducing one
generally increases another; there is no free lunch.
\end{concept}

\begin{analogy}
Imagine three groups of archers at a target. The high-bias group clusters
all arrows together, but consistently far from the bullseye. The
high-variance group spreads arrows widely in every direction. The
well-calibrated group clusters near the center. Good modeling is the
third group: consistent and accurate.
\end{analogy}

\begin{equation*}
  E\!\left[(y - \hat{f}(x))^2\right]
  = \Bigl(E[\hat{f}(x)] - f(x)\Bigr)^2
    + \mathrm{Var}[\hat{f}(x)]
    + \sigma^2
\end{equation*}

The expectation is taken over the randomness in both the training data
and the label noise $\sigma^2$.

\begin{mentalmodel}
The expectation in the decomposition is over hypothetical re-draws of
entire training sets, not over individual noisy labels. Bias and variance
describe properties of the learning algorithm across many possible
datasets, not the behavior of one fitted model on one dataset. That
distinction is why you cannot estimate bias from a single run.
\end{mentalmodel}

\vspace{10pt}\hrule height 0.4pt\vspace{14pt}

% ---------------------------------------------------------------------------
\section*{Regularization and the Art of Not Overfitting}
\addcontentsline{toc}{section}{Regularization and the Art of Not Overfitting}
% ---------------------------------------------------------------------------

Regularization often feels like an arbitrary penalty added to the loss.
It is not arbitrary at all: it is a prior belief about the size of
coefficients, expressed in the language of optimization. Where do your
priors come from, and what do they imply about the penalties you add?

\begin{concept}
L2 regularization is MAP estimation under a Gaussian prior on the
coefficients; L1 is MAP under a Laplace prior. The penalty strength
$\lambda = \sigma^2/\tau^2$ is the ratio of noise variance to prior
variance, exactly analogous to the complement $1-Z$ in Buhlmann
credibility. Regularization is Bayesian shrinkage written in
optimization notation.
\end{concept}

\begin{analogy}
Think of regularization as a complexity budget. L2 charges quadratically:
spreading weight across many features is affordable because each small
weight costs little. L1 charges linearly: a single large weight is
as expensive as many small ones, so the budget pressure pushes toward
sparse solutions where most weights land at exactly zero.
\end{analogy}

The log-posterior under a Gaussian prior recovers Ridge regression:

\begin{equation*}
  \log\pi(\beta \mid y)
  \propto
  -\frac{1}{2\sigma^2}\|y - X\beta\|_2^2
  - \frac{1}{2\tau^2}\|\beta\|_2^2
\end{equation*}

Setting $\lambda = \sigma^2/\tau^2$ gives the standard Ridge objective.

\begin{mentalmodel}
Ridge regression is multivariate Buhlmann credibility: each coefficient is
shrunk toward zero by the ratio of noise variance to prior variance, just
as an individual claim mean is blended toward the portfolio mean in
proportion to how much data supports it. The math is identical; only the
vocabulary differs.
\end{mentalmodel}

\vspace{10pt}\hrule height 0.4pt\vspace{14pt}

% ---------------------------------------------------------------------------
\section*{Cross-Validation Strategies}
\addcontentsline{toc}{section}{Cross-Validation Strategies}
% ---------------------------------------------------------------------------

Cross-validation is easy to run and easy to misinterpret. The metric it
returns is only as trustworthy as the splitting scheme used to produce it.
Have you checked whether your splits respect the actual structure of your
data?

\begin{concept}
Cross-validation averages held-out performance estimates to get a low-variance
proxy for generalization error. The critical requirement is that the splitting
scheme must reflect how the model will be used in production. A split that
does not match the deployment structure produces an estimate that is
confidently wrong, which is worse than no estimate at all.
\end{concept}

\begin{analogy}
Each fold is a dress rehearsal under one specific scenario. The mean score
across folds is the expected performance; the spread across folds tells you
how fragile the model is to which slice of data happens to be held out.
A low mean with high spread signals that performance is driven by a few
easy subsets, not by general learning.
\end{analogy}

A key theoretical result due to Bengio and Grandvalet (2004) is that no
general unbiased estimator of the variance of k-fold CV error exists.
The fold standard deviation should be treated as a stability indicator,
not as a classical confidence interval.

\begin{mentalmodel}
For time series, the arrow of time is inviolable. Purged and embargoed
cross-validation adds a gap at fold boundaries to prevent features or
labels computed near the boundary of one fold from bleeding information
into the adjacent fold. Skipping this step produces optimistically biased
estimates that collapse the moment the model is deployed.
\end{mentalmodel}

\vspace{10pt}\hrule height 0.4pt\vspace{14pt}

% ---------------------------------------------------------------------------
\section*{Data Leakage Patterns}
\addcontentsline{toc}{section}{Data Leakage Patterns}
% ---------------------------------------------------------------------------

Data leakage is the most reliably damaging mistake in applied machine
learning, and it is easy to commit without noticing. The diagnostic
question is deceptively simple: at the moment of prediction in production,
exactly what information is available?

\begin{concept}
Anything used during training that would not be available at prediction
time is leakage. Leakage causes offline metrics to overstate real-world
performance, sometimes dramatically. The model has learned to exploit an
information shortcut that disappears the moment it is deployed.
\end{concept}

\begin{analogy}
Imagine an exam where the answer key was accidentally included in the exam
packet. The student scores perfectly, but learned nothing about the subject.
Offline AUC looks excellent because the model memorized the leak, not the
underlying pattern. The course correction feels brutal precisely because the
misleading signal was so convincing.
\end{analogy}

The canonical fix is a scikit-learn \texttt{Pipeline} that fits all
transformers only on training folds, never on the validation fold.
Point-in-time feature stores enforce temporal cutoffs at the row level,
ensuring that no feature for a given row uses information from after the
prediction timestamp.

\begin{mentalmodel}
A too-good-to-be-true offline metric, such as AUC 0.99 on a genuinely
hard real-world problem, is the primary warning sign. A single feature
that explains almost everything is the secondary one. Either observation
should trigger an immediate audit of the data pipeline, not a celebration.
\end{mentalmodel}

\vspace{10pt}\hrule height 0.4pt\vspace{14pt}

% ---------------------------------------------------------------------------
\section*{Feature Engineering}
\addcontentsline{toc}{section}{Feature Engineering}
% ---------------------------------------------------------------------------

Raw columns rarely tell the whole story. Feature engineering is the
systematic process of asking: what does an expert know that the data
does not yet express? The answer to that question is often the difference
between a mediocre model and a great one.

\begin{concept}
Features encode domain knowledge in a form the model can exploit. The
question is always whether the transformation reveals structure that the
raw data obscures. Target encoding and actuarial credibility are the same
estimator: both shrink a group mean toward a global mean by an amount
proportional to how much data the group contains.
\end{concept}

\begin{analogy}
The smoothed target encoding for category $c$ is
$\hat{y}_c = Z_c\bar{y}_c + (1-Z_c)\bar{y}$, where
$Z_c = n_c/(n_c + k)$. This is Buhlmann's formula with
$k = \sigma^2_\varepsilon/\sigma^2_\mu$: the credibility weight grows as
the group provides more observations and shrinks when the group is noisy
relative to the portfolio.
\end{analogy}

Weight of Evidence and Information Value are classical credit-scoring
features that generalize to any binary classification problem:

\begin{equation*}
  \mathrm{WoE}_c = \log\frac{P_c}{N_c},
  \qquad
  \mathrm{IV} = \sum_c (P_c - N_c)\cdot\mathrm{WoE}_c
\end{equation*}

where $P_c$ and $N_c$ are the shares of positives and negatives
falling in category $c$.

\begin{mentalmodel}
WoE converts each category to a log-odds score: positive values signal
enrichment for positive outcomes, negative values signal depletion. IV
sums these per-category log-odds signals, weighted by how common each
category is, into a single number that summarizes how much discriminatory
power a feature carries across all its levels.
\end{mentalmodel}

\vspace{10pt}\hrule height 0.4pt\vspace{14pt}

% =============================================================================
\part*{Part III: Inference, Causality, and Uncertainty}
\addcontentsline{toc}{part}{Part III: Inference, Causality, and Uncertainty}
% =============================================================================

% ---------------------------------------------------------------------------
\section*{Randomized Experiments and A/B Testing}
\addcontentsline{toc}{section}{Randomized Experiments and A/B Testing}
% ---------------------------------------------------------------------------

Correlation shows up everywhere in product data. Causation is rarer, and
only one method establishes it by construction. What is special about
randomization, and why does it justify causal claims that no amount of
observational data can?

\begin{concept}
Randomization breaks the dependence between treatment assignment and every
possible confounder, observed or unobserved. As a result, the observed
difference in outcomes between treatment and control is an unbiased
estimate of the average treatment effect. No other study design achieves
this without additional assumptions.
\end{concept}

\begin{analogy}
An A/B test is a controlled scientific experiment embedded in a live
product. The discipline of a good experiment is mostly pre-experiment
work: pre-register the hypothesis, metric, minimum detectable effect,
required duration, and decision rule before the test starts. Any
analysis conducted after peeking at results is exploratory, not
confirmatory.
\end{analogy}

The required sample size per arm for a two-sample test is:

\begin{equation*}
  n_{\text{per arm}}
  = \frac{2\,(z_{1-\alpha/2} + z_{1-\beta})^2\,\sigma^2}{\delta^2}
\end{equation*}

where $\delta$ is the minimum detectable effect, $\sigma^2$ is the
outcome variance, and $z_{1-\alpha/2}$, $z_{1-\beta}$ are the normal
quantiles for the chosen Type I and Type II error rates.

\begin{mentalmodel}
CUPED variance reduction is essentially free. Residualizing the outcome
on a pre-experiment covariate reduces variance by a factor of
$1 - \rho^2$. At $\rho = 0.7$, a 20{,}000-user experiment shrinks to
roughly 10{,}200 users for the same power, cutting the required runtime
nearly in half with no additional cost.
\end{mentalmodel}

\vspace{10pt}\hrule height 0.4pt\vspace{14pt}

% ---------------------------------------------------------------------------
\section*{Observational Causal Inference}
\addcontentsline{toc}{section}{Observational Causal Inference}
% ---------------------------------------------------------------------------

We cannot always run experiments. Treatment may be unethical to withhold,
impossible to control, or already baked into historical data. Observational
causal inference asks: under what assumptions can we recover causal effects
from data that was not randomized?

\begin{concept}
Potential outcomes $Y(1)$ and $Y(0)$ are never simultaneously observed for
the same unit. Causal identification requires assumptions: parallel trends
for difference-in-differences, ignorability for propensity score methods,
or an exclusion restriction for instrumental variables. The assumptions are
the argument; the method is just the arithmetic that follows from them.
\end{concept}

\begin{analogy}
A directed acyclic graph is the causal map you draw before touching data.
Nodes are variables, arrows are direct causal influences. Backdoor paths
from treatment to outcome through confounders are what create spurious
associations. The method you choose should be the one that blocks those
paths given what you can observe and control.
\end{analogy}

The difference-in-differences estimator is:

\begin{equation*}
  \hat\tau_{\mathrm{DiD}}
  = \bigl(\bar{Y}_{T,\text{post}} - \bar{Y}_{T,\text{pre}}\bigr)
  - \bigl(\bar{Y}_{C,\text{post}} - \bar{Y}_{C,\text{pre}}\bigr)
\end{equation*}

Identification requires that, absent treatment, the treatment group would
have followed the same trend as the control group (parallel trends).

\begin{mentalmodel}
CausalImpact uses a Bayesian structural time series model to construct a
synthetic counterfactual from the pre-period and a set of donor series.
This is DiD for a single treated time series: the pre-period model is the
counterfactual; the post-period gap between actuals and the synthetic
control is the estimated treatment effect.
\end{mentalmodel}

\vspace{10pt}\hrule height 0.4pt\vspace{14pt}

% ---------------------------------------------------------------------------
\section*{Calibration and Conformal Prediction}
\addcontentsline{toc}{section}{Calibration and Conformal Prediction}
% ---------------------------------------------------------------------------

A model that ranks cases perfectly can still be wrong about its
probabilities. Calibration and discrimination are independent properties,
and real decisions often require both. Would you trust a model that says
``90\% confident'' when it is right only 60\% of the time?

\begin{concept}
A model is calibrated if, among all examples assigned probability $p$,
a fraction $p$ are actually positive. Calibration is a global property
of the probability outputs, not of the ranking. A model can have excellent
AUC and terrible calibration, or the reverse. For risk pricing and
reserve setting, calibration is not optional.
\end{concept}

\begin{analogy}
A weather forecaster who says ``70\% chance of rain'' should be right
on 70\% of the days they make that call. The reliability diagram plots
actual frequency against predicted probability across confidence bins;
a perfectly calibrated model falls on the diagonal. ECE summarizes the
average gap into one number:
\end{analogy}

\begin{equation*}
  \mathrm{ECE}
  = \sum_{m} \frac{|B_m|}{n}
    \left|\mathrm{acc}(B_m) - \mathrm{conf}(B_m)\right|
\end{equation*}

where $B_m$ is a confidence bin, $\mathrm{acc}(B_m)$ is the fraction
positive in that bin, and $\mathrm{conf}(B_m)$ is the mean predicted
probability.

\begin{mentalmodel}
Conformal prediction provides the following finite-sample guarantee under
exchangeability alone, with no model assumptions:
$P(Y_{n+1} \in C(X_{n+1})) \geq 1-\alpha$.
No distributional assumption is needed; only exchangeability. That makes
conformal prediction sets the most defensible form of uncertainty
quantification for a regulatory or audit context.
\end{mentalmodel}

\vspace{10pt}\hrule height 0.4pt\vspace{14pt}

% =============================================================================
\part*{Part IV: Interpretability, Fairness, and Governance}
\addcontentsline{toc}{part}{Part IV: Interpretability, Fairness, and Governance}
% =============================================================================

% ---------------------------------------------------------------------------
\section*{Model Interpretability with SHAP and Friends}
\addcontentsline{toc}{section}{Model Interpretability with SHAP and Friends}
% ---------------------------------------------------------------------------

Knowing a model is accurate is not the same as knowing why it is accurate.
For decisions that affect people, the explanation matters as much as the
prediction. SHAP values offer one mathematically principled way to
distribute credit among features.

\begin{concept}
SHAP values are the unique attribution satisfying four Shapley axioms:
efficiency (contributions sum to the prediction minus the baseline),
symmetry (equal contributors get equal credit), null player (features that
change nothing get zero), and additivity (SHAP values of combined models
sum correctly). Efficiency is the one most practitioners care about in
production: every gap between prediction and baseline is fully explained.
\end{concept}

\begin{analogy}
Imagine a group of colleagues who collaborated on a project with a
measurable outcome. Each colleague's fair share is their average marginal
contribution: how much better the outcome gets when they are added to every
possible subset of the team, averaged over all orderings. SHAP applies
exactly that fairness principle to features in a model.
\end{analogy}

\begin{equation*}
  \phi_j = \sum_{S \subseteq F\setminus\{j\}}
  \frac{|S|!\,(|F|-|S|-1)!}{|F|!}
  \Bigl[v(S\cup\{j\}) - v(S)\Bigr]
\end{equation*}

where $v(S)$ is the expected model prediction when only features in $S$
are known and $F$ is the full feature set.

\begin{mentalmodel}
Conditional SHAP and interventional SHAP give different answers when
features are correlated. Conditional SHAP respects the data distribution;
interventional SHAP treats features as independent. For regulatory
adverse-action letters, document which value function you used and why,
because the choice changes which features receive credit.
\end{mentalmodel}

\vspace{10pt}\hrule height 0.4pt\vspace{14pt}

% ---------------------------------------------------------------------------
\section*{Fairness and Responsible AI}
\addcontentsline{toc}{section}{Fairness and Responsible AI}
% ---------------------------------------------------------------------------

There is no single, universally agreed definition of fairness. That is not
a philosophical footnote; it is a practical constraint that forces a
decision before any model is deployed. Which definition you target, and why,
is a normative choice that belongs to the humans involved, not to the model.

\begin{concept}
Demographic parity, equal opportunity, equalized odds, and predictive
parity are mutually incompatible whenever base rates differ across groups.
Satisfying one definition mathematically precludes satisfying the others.
There is no technical fix; the tension is fundamental, and the resolution
requires a policy decision about which errors are more acceptable.
\end{concept}

\begin{analogy}
Intersectional analysis reveals gaps that single-attribute analysis hides.
A model can satisfy the 80\% rule on gender and separately on ethnicity,
while producing substantially worse outcomes for a specific intersection
such as indigenous women, than any one-dimensional audit would detect.
Fairness audits that ignore intersections are incomplete by construction.
\end{analogy}

Differential privacy provides a formal bound on the information a model
reveals about any individual record:

\begin{equation*}
  P(M(D)\in S) \leq e^\epsilon\, P(M(D')\in S) + \delta
\end{equation*}

where $D$ and $D'$ differ in exactly one record. Smaller $\epsilon$ means
stronger privacy guarantees; $\delta$ is the probability of a catastrophic
failure of the bound.

\begin{mentalmodel}
Fairness is not a setting to enable; it is a conversation to conduct.
Declare which definition you target and justify it. Measure the gap per
group before and after any intervention. Document the accuracy-fairness
tradeoff explicitly. The documentation is evidence that the decision was
made deliberately rather than by default.
\end{mentalmodel}

\vspace{10pt}\hrule height 0.4pt\vspace{14pt}

% =============================================================================
\part*{Part V: Modern Technique Families}
\addcontentsline{toc}{part}{Part V: Modern Technique Families}
% =============================================================================

% ---------------------------------------------------------------------------
\section*{Deep Learning Fundamentals}
\addcontentsline{toc}{section}{Deep Learning Fundamentals}
% ---------------------------------------------------------------------------

Neural networks looked like they would require enormous amounts of
hand-crafted gradient code until automatic differentiation changed
everything. What made deep learning suddenly tractable was not just more
data or faster chips; it was backpropagation built into software that
could handle arbitrary computation graphs.

\begin{concept}
Backpropagation is the chain rule of calculus applied systematically to
a computational graph. Modern autograd frameworks build the graph during
the forward pass and traverse it backward automatically. The user writes
the forward computation; the framework computes all gradients without any
user-written derivative code.
\end{concept}

\begin{analogy}
A ResNet can be understood as an exponential ensemble of shallow paths.
Each residual block provides an identity shortcut: the block's output
can be skipped entirely or added to the learned transformation, giving
$2^L$ implicit subnetworks across $L$ blocks. Gradients always have a
direct channel back through the identity paths, which is why very deep
ResNets remain trainable.
\end{analogy}

LoRA (Low-Rank Adaptation) fine-tunes a large model by adding a low-rank
update to each weight matrix:

\begin{equation*}
  W' = W + \Delta W,
  \quad
  \Delta W = B A,
  \quad
  A \in \mathbb{R}^{r \times n},\;
  B \in \mathbb{R}^{m \times r},\;
  r \ll \min(m,n)
\end{equation*}

Only $r(m+n)$ parameters are trained instead of $mn$, reducing memory
and compute by orders of magnitude for large weight matrices.

\begin{mentalmodel}
Training a deep network is like teaching a chain of translators: each
layer translates the representation from the previous layer into something
more useful for the next. Residual connections add a direct phone line so
that error signals can reach every layer without being distorted by
passing through every intermediate translation step.
\end{mentalmodel}

\vspace{10pt}\hrule height 0.4pt\vspace{14pt}

% ---------------------------------------------------------------------------
\section*{LLM Foundations for Data Scientists}
\addcontentsline{toc}{section}{LLM Foundations for Data Scientists}
% ---------------------------------------------------------------------------

Large language models are not magic; they are transformers trained on
next-token prediction at a scale that produces emergent capabilities.
Understanding their training pipeline helps you decide when to use them,
how to adapt them, and where they will fail.

\begin{concept}
LLMs are trained in stages: pretraining on next-token prediction over vast
corpora, supervised fine-tuning on instruction-response pairs, alignment
via RLHF or DPO to match human preferences, and domain adaptation via
low-rank methods such as LoRA. Each stage adds a layer of behavior on top
of the statistical foundation built by pretraining.
\end{concept}

\begin{analogy}
Speculative decoding is free speed. A small draft model proposes $k$ tokens
cheaply; the large target model verifies all $k$ in a single forward pass,
which costs roughly the same as generating one token. Accepted tokens
commit; the first rejection rolls back. The result is a 2x to 3x latency
reduction with no change in output quality.
\end{analogy}

Scaled dot-product attention is the core computational primitive:

\begin{equation*}
  \mathrm{Attention}(Q,K,V)
  = \mathrm{softmax}\!\left(\frac{QK^\top}{\sqrt{d_k}}\right)V
\end{equation*}

The $\sqrt{d_k}$ scaling prevents dot products from growing so large that
they push the softmax into a saturated regime where gradients vanish.

\begin{mentalmodel}
Attention is a soft database lookup. The query says ``what am I looking
for''; the keys say ``what each stored entry is about''; the values say
``what each entry actually contains.'' The softmax computes a weighted
blend: entries whose keys match the query closely contribute most to the
output.
\end{mentalmodel}

\vspace{10pt}\hrule height 0.4pt\vspace{14pt}

% ---------------------------------------------------------------------------
\section*{Computer Vision Foundations}
\addcontentsline{toc}{section}{Computer Vision Foundations}
% ---------------------------------------------------------------------------

Image pixels are not meaningful in isolation. What matters is the spatial
relationship between pixels: an edge, a texture, a shape. Convolutional
layers exploit this structure explicitly, and the resulting visual
hierarchy turns out to transfer remarkably well across domains.

\begin{concept}
Convolutional layers are local, weight-shared, and parameter-efficient.
Each filter learns to detect a particular pattern anywhere in the image.
As the network deepens, filters at higher layers respond to increasingly
abstract patterns: edges, then textures, then parts, then objects.
This hierarchy is the source of both the efficiency and the transferability
of convolutional networks.
\end{concept}

\begin{analogy}
CLIP's zero-shot classification illustrates how vision and language can
be aligned without task-specific training. Encode each candidate label
as ``a photo of a [class]'' and encode the input image. Predicted class
is the label with highest cosine similarity. No labeled examples of the
new class are required; the shared embedding space does the work.
\end{analogy}

Diffusion models learn to generate images by reversing a forward noising
process. One step of the forward process is:

\begin{equation*}
  q(x_t \mid x_{t-1})
  = \mathcal{N}\!\bigl(x_t;\;
    \sqrt{1-\beta_t}\,x_{t-1},\;
    \beta_t I\bigr)
\end{equation*}

After $T$ steps, the image becomes approximately isotropic Gaussian noise;
the network learns to denoise backward from $T$ to $0$.

\begin{mentalmodel}
Transfer learning works because the visual hierarchy from edges to objects
is domain-general. The features learned on ImageNet are not about
ImageNet; they are about the geometry of natural images. Only the final
mapping from abstract visual concepts to your specific class labels needs
to be retrained.
\end{mentalmodel}

\vspace{10pt}\hrule height 0.4pt\vspace{14pt}

% ---------------------------------------------------------------------------
\section*{Time Series Foundations}
\addcontentsline{toc}{section}{Time Series Foundations}
% ---------------------------------------------------------------------------

Time series data has one feature that distinguishes it from every other
data type: the order of rows encodes information. Shuffle a tabular dataset
and you still have the same dataset; shuffle a time series and you have
destroyed it. That ordering constraint permeates modeling, evaluation,
and feature engineering.

\begin{concept}
Time series observations are sequentially ordered; swapping rows changes
the data. This violates the IID assumption underlying most off-the-shelf
ML methods and requires specialized evaluation (forward-chaining CV),
modeling (ARIMA, GARCH, neural sequence models), and feature engineering
(lag features, rolling statistics, Fourier terms).
\end{concept}

\begin{analogy}
GARCH is a weather model for volatility. Calm days tend to be followed by
calm days; stormy days tend to be followed by stormy days. The model tracks
the running mood of the market and adjusts the expected variance for the
next period based on recent shocks and recent conditional variances.
\end{analogy}

The GARCH$(p,q)$ conditional variance equation:

\begin{equation*}
  \sigma_t^2
  = \alpha_0
  + \sum_{i=1}^{q} \alpha_i\,\varepsilon_{t-i}^2
  + \sum_{j=1}^{p} \beta_j\,\sigma_{t-j}^2
\end{equation*}

Stationarity of the variance process requires
$\sum_{i}\alpha_i + \sum_{j}\beta_j < 1$.

\begin{mentalmodel}
Hierarchical forecasting is the problem: a multi-level structure such as
nation, region, store, SKU must produce forecasts that add up consistently.
Forecast reconciliation is the family of solutions. The MinT method
chooses reconciliation weights that minimize the total forecast-error
variance across all levels of the hierarchy.
\end{mentalmodel}

\vspace{10pt}\hrule height 0.4pt\vspace{14pt}

% ---------------------------------------------------------------------------
\section*{Recommender System Foundations}
\addcontentsline{toc}{section}{Recommender System Foundations}
% ---------------------------------------------------------------------------

Recommender systems face a data problem that most classification tasks
do not: the absence of an interaction does not mean the user dislikes the
item. It may simply mean the user never saw it. How you handle that
asymmetry between observed and unobserved interactions shapes everything
downstream.

\begin{concept}
Implicit feedback such as clicks and purchases is abundant but ambiguous.
Unobserved items are not negatives; they are missing. Pairwise losses
that require only relative ordering between observed and unobserved items
are better suited to this data structure than pointwise MSE, which treats
every unobserved item as a zero-score target.
\end{concept}

\begin{analogy}
A two-stage recommender system resembles airport security. Retrieval,
typically approximate nearest neighbor search in an embedding space, is
the fast first pass that narrows millions of candidates to a shortlist of
hundreds. Ranking is the thorough second pass that scores each candidate
carefully using a richer feature set.
\end{analogy}

Bayesian Personalized Ranking optimizes for the probability that observed
items rank above unobserved items:

\begin{equation*}
  \mathcal{L}_{\text{BPR}}
  = \sum_{(u,i,j)\in D_S}
    \log\sigma(\hat{x}_{ui} - \hat{x}_{uj})
    - \lambda\|\Theta\|^2
\end{equation*}

where $i$ is an observed interaction, $j$ is an unobserved item, and
$\sigma$ is the sigmoid function.

\begin{mentalmodel}
Off-policy evaluation via inverse propensity scoring is the same
reweighting technique as in observational causal inference. Logged
interactions were generated under a historical recommendation policy;
IPS reweights them to estimate what a new policy would have produced,
allowing you to evaluate a new ranker without running a live experiment.
\end{mentalmodel}

\vspace{10pt}\hrule height 0.4pt\vspace{14pt}

% =============================================================================
\part*{Part VI: Production and Infrastructure}
\addcontentsline{toc}{part}{Part VI: Production and Infrastructure}
% =============================================================================

% ---------------------------------------------------------------------------
\section*{MLOps Architecture Patterns}
\addcontentsline{toc}{section}{MLOps Architecture Patterns}
% ---------------------------------------------------------------------------

A model that works in a notebook and a model that works in production are
different objects. The gap between them is the subject of MLOps: not just
deployment, but monitoring, retraining triggers, data versioning, and
feature consistency. What breaks first when a model moves from a
researcher's laptop to a production API?

\begin{concept}
Training-serving skew is the most common silent production failure.
Features computed differently at training time versus inference time cause
the model to receive inputs outside its training distribution without any
visible error. A feature store with point-in-time correctness is the
primary defense: it enforces that training and serving retrieve features
using identical logic.
\end{concept}

\begin{analogy}
Shadow deployment is the safest way to evaluate a candidate model.
Every production request is duplicated; the candidate model receives a
copy and produces a prediction that is logged but never acted on. Users
see no change; you accumulate full-fidelity production signal to compare
the candidate against the current model before any traffic is shifted.
\end{analogy}

Population Stability Index measures distribution shift between training
and production:

\begin{equation*}
  \mathrm{PSI}
  = \sum_i (p_i - q_i)\log\frac{p_i}{q_i}
\end{equation*}

where $p_i$ is the training-time proportion and $q_i$ is the current
production proportion in bin $i$. PSI above 0.2 is the conventional
threshold for investigation.

\begin{mentalmodel}
MLOps maturity is a step function, not a slope. Maturity 0 is deploying
a notebook manually. Maturity 1 is an automated training and deployment
pipeline triggered by new data. Maturity 2 is CI/CD for the model itself:
code changes trigger tests, performance regressions trigger rollbacks.
Each level is a qualitative improvement in production reliability.
\end{mentalmodel}

\vspace{10pt}\hrule height 0.4pt\vspace{14pt}

% ---------------------------------------------------------------------------
\section*{Cloud Deployment for ML Services}
\addcontentsline{toc}{section}{Cloud Deployment for ML Services}
% ---------------------------------------------------------------------------

A trained model sitting in a file has zero business value. Converting it
into a service that responds to requests, scales with demand, and recovers
from failures is the deployment problem. What does a minimal but
production-grade deployment look like, and what does each component
actually protect against?

\begin{concept}
Cloud deployment converts a trained artifact into a service with
autoscaling, monitoring, and rollback. Without these capabilities, the
model cannot serve users reliably, cannot adapt to load changes, and
cannot recover from errors without manual intervention. Deployment is not
the end of the ML workflow; it is where the workflow begins to matter.
\end{concept}

\begin{analogy}
A multi-service request taking 2.5 seconds may spend 200 ms inside the
model and 2.3 seconds waiting for a misconfigured database view. Without
distributed tracing you blame the model; with OpenTelemetry traces you
see the actual bottleneck in the waterfall and fix it in minutes rather
than days of speculation.
\end{analogy}

A minimal production deployment can be thought of as five concentric
reliability rings: the Docker image (reproducibility), Cloud Run
(autoscaling and managed infrastructure), GitHub Actions CI/CD
(automated testing and deployment), a monitoring dashboard
(observability), and budget alerts (cost control). Each ring catches a
different class of failure that the inner rings cannot prevent.

\begin{mentalmodel}
Cold-start latency is the main user-experience risk in serverless
deployment. Mitigations layer: lazy model loading defers expensive work
until the first request, startup probes prevent traffic from reaching
an unready instance, provisioned concurrency keeps warm instances ready,
and model-server warmup caches predictions for common inputs. Together,
these keep p99 latency within SLA even when instances spin up from zero.
\end{mentalmodel}

\vspace{10pt}\hrule height 0.4pt\vspace{14pt}

% ---------------------------------------------------------------------------
\section*{SQL for Data Scientists}
\addcontentsline{toc}{section}{SQL for Data Scientists}
% ---------------------------------------------------------------------------

SQL is four decades old and remains the most widely used language for
data manipulation in production. Its staying power comes from a design
choice that is easy to underestimate: SQL is declarative. You describe
the result; the query planner decides how to compute it. Understanding
the planner is what separates a query that runs in seconds from one that
runs for hours.

\begin{concept}
SQL is declarative: you specify the shape of the result, not the steps
to compute it. The query planner translates your intent into a physical
execution plan. Your job is to help the planner by providing accurate
schema information, appropriate indexes, and query shapes that expose
the structure the planner can exploit, such as filter pushdown and
partition pruning.
\end{concept}

\begin{analogy}
The ROWS versus RANGE distinction in window functions is invisible until
there are tied values. Consider three orders placed on the same date.
ROWS treats each row as a separate point in the frame, counting them one
by one. RANGE treats all rows with the same ORDER BY value as a peer
group and includes all of them simultaneously. The results diverge
exactly when ties exist and exactly where correct lifetime-value
calculation matters most.
\end{analogy}

The standard running-total window function uses a ROWS frame:
\texttt{SUM(revenue) OVER (PARTITION BY customer\_id ORDER BY order\_date
ROWS BETWEEN UNBOUNDED PRECEDING AND CURRENT ROW)}. This computes the
running lifetime value per customer, adding each order one row at a time
and never including future rows in the cumulative total.

\begin{mentalmodel}
\texttt{EXPLAIN ANALYZE} shows the plan the query planner actually chose
and the real row counts at each step. A sequential scan on a large table
indicates a missing index. A row-count estimate that is orders of magnitude
from the actual count indicates stale statistics; running
\texttt{ANALYZE} on the table lets the planner rebuild its estimates and
often halves execution time without changing the query.
\end{mentalmodel}

\vspace{10pt}\hrule height 0.4pt\vspace{14pt}

% =============================================================================
\section*{Cheat Sheet}
\addcontentsline{toc}{section}{Cheat Sheet}
% =============================================================================

\begin{center}
\small
\begin{tabular}{@{}p{4.8cm} p{5.2cm} p{4.5cm}@{}}
\toprule
\textbf{Concept} & \textbf{One-Line Model} & \textbf{Picture} \\
\midrule
Probability \& MLE
  & $\sqrt{n}(\hat\theta-\theta_0)\to\mathcal{N}(0,I^{-1})$
  & Sharpness of likelihood peak \\
\addlinespace
Optimization \& Loss
  & AdamW: adaptive rates with decoupled weight decay
  & Condition number = bowl aspect ratio \\
\addlinespace
Glossary
  & Precision in language is precision in thinking
  & Shared glossary prevents week-long wrong models \\
\addlinespace
Bias-Variance
  & $E[(y-\hat f)^2] = \text{Bias}^2 + \text{Var} + \sigma^2$
  & Three archers: scattered, biased, centered \\
\addlinespace
Regularization
  & L2 = Gaussian MAP; L1 = Laplace MAP; $\lambda=\sigma^2/\tau^2$
  & Complexity budget charged quadratically or linearly \\
\addlinespace
Cross-Validation
  & Mean score = performance; spread = fragility
  & Dress rehearsal under each scenario \\
\addlinespace
Data Leakage
  & Available at prediction time? If not, it is leakage.
  & Answer key in the exam packet \\
\addlinespace
Feature Engineering
  & WoE, target encoding = Buhlmann credibility
  & Expert knowledge encoded as a number \\
\addlinespace
A/B Testing
  & $n = 2(z_{\alpha/2}+z_\beta)^2\sigma^2/\delta^2$
  & Pre-registered scientific experiment in a product \\
\addlinespace
Causal Inference
  & DiD = double difference; DAG = causal map
  & Block backdoor paths before touching data \\
\addlinespace
Calibration
  & ECE = mean gap between confidence and accuracy
  & Weather forecaster: 70\% call should hit 70\% of the time \\
\addlinespace
SHAP
  & $\phi_j$ = average marginal contribution over all subsets
  & Team credit split fairly by marginal contribution \\
\addlinespace
Fairness
  & No single definition; declare which one and measure gaps
  & Intersections hide what single attributes conceal \\
\addlinespace
Deep Learning
  & Backprop = chain rule on computation graph
  & Chain of translators with direct phone lines (residuals) \\
\addlinespace
LLM Foundations
  & Attention = soft database lookup; LoRA = low-rank update
  & Pretraining + SFT + alignment + domain adaptation \\
\addlinespace
Computer Vision
  & Conv layers: local, shared, hierarchical
  & Visual hierarchy transfers; only the head needs retraining \\
\addlinespace
Time Series
  & GARCH$(p,q)$: $\sigma^2_t=\alpha_0+\sum\alpha_i\epsilon^2_{t-i}+\sum\beta_j\sigma^2_{t-j}$
  & Volatility clusters; order of rows encodes information \\
\addlinespace
Recommender Systems
  & BPR: observed items should rank above unobserved
  & Two-stage: retrieval (fast) then ranking (careful) \\
\addlinespace
MLOps
  & PSI $>$ 0.2 triggers drift investigation
  & Maturity 0: notebook; 1: pipeline; 2: CI/CD for model \\
\addlinespace
Cloud Deployment
  & Docker + Cloud Run + CI/CD + monitoring + budget alerts
  & OpenTelemetry traces reveal the real bottleneck \\
\addlinespace
SQL
  & Declarative: describe result; planner chooses method
  & EXPLAIN ANALYZE: find the sequential scan, fix the index \\
\addlinespace
Reading Guide
  & Purple = what; amber = feels like; green = remember
  & Three box colors, fixed meaning, no title needed \\
\bottomrule
\end{tabular}
\end{center}

\end{document}
